\section{练习}

\begin{question}\begin{drill}[螺线管内场变化与不等半圆]\label{em2:drill:1}
长螺线管截面半径 $a$，电流为 $I_0$ 时内场为 $B_0$。电流以速率 $\beta>0$ 增加；管外同轴圆环由两个半圆构成，正向 $P\to Q$ 的半圆电阻为 $R_1$，另一半为 $R_2$。选正法线使总电动势为正，求环流和 $V_P-V_Q$。
\end{drill}
\fig{fig:em2:drill:1:s}{外环包围整个螺线管截面；$a$ 是磁区半径，外环半径另记为 $\rho>a$}{\begin{tikzpicture}

\fill[blue!8](0,0)circle(.9);\draw[world](0,0)circle(.9)circle(1.7);\draw[dim](0,0)--(.9,0)node[midway,below]{$a$};\draw[dim](0,0)--({1.7*cos(140)},{1.7*sin(140)})node[midway,left]{$\rho$};\node at(0,.4){$\otimes\dot B$};\node[right]at(1.7,0){$P$};\node[left]at(-1.7,0){$Q$};\node at(0,2.3){$R_1:\ P\to Q$};\node at(0,-2.3){$R_2:\ Q\to P$};\draw[->,blue](1.7,0)arc(0:75:1.7);

\end{tikzpicture}}
\end{question}
\answerlink{em2:drill:1}
\begin{question}\begin{drill}[绕中线转动的等边三角形]\label{em2:drill:2}
边长 $l$、总电阻 $R$ 的均匀等边三角形 $abd$ 绕中线 $ac$ 转动，$c$ 为底边中点。此刻线框在 $xy$ 平面，$a$ 在 $+y$ 轴，$b$ 在 $+x$ 一侧。磁场沿 $+x$，角速度沿 $+y$，大小为 $B,\omega$。求总电动势及 $U_{ab},U_{ac}$。
\end{drill}
\fig{fig:em2:drill:2:s}{题设：绕中线转动的等边三角形}{\begin{tikzpicture}[scale=.88]\draw[world](0,2.6)node[above]{$a$}--(1.5,0)node[below right]{$b$}--(-1.5,0)node[below left]{$d$}--cycle;\node[below]at(0,0){$c$};\draw[guide](0,-.4)--(0,3.2)node[above]{$\omega,+y$};\draw[flow](2,1)--(3,1)node[right]{$B,+x$};\end{tikzpicture}}
\end{question}
\answerlink{em2:drill:2}
\begin{question}\begin{drill}[改变引线面积造成的读数变化]\label{em2:drill:3}
固定的两端点 $A,B$ 连接一只内阻足够大的电压表。两次引线走向分别为 $g_1,g_2$，都按 $A\to B$ 取向。沿 $g_1$ 前进、$g_2$ 返回构成面积 $S$ 的矩形，正法线向纸外。该窄条内 $\dot B=\beta$ 向纸外。求两次有向读数之差及误差容许 $u_0$ 对面积的要求。
\end{drill}
\fig{fig:em2:drill:3:s}{题设：改变引线面积造成的读数变化}{\begin{tikzpicture}[scale=.88]\draw[world](0,0)node[left]{$A$}--(4,0)node[right]{$B$};\draw[guide](0,0)--(0,1.5)--(4,1.5)--(4,0);\node[below]at(2,0){$g_1$};\node[above]at(2,1.5){$g_2$};\node at(2,.8){$\odot\dot B=\beta,\ S$};\end{tikzpicture}}
\end{question}
\answerlink{em2:drill:3}
\begin{question}\begin{drill}[极性相反的两电容并接]\label{em2:drill:4}
电容 $C,2C$ 的上极板初电压分别为 $+U,-U$。用细电阻将两上极板相连，下极板也相连，不接外源，最终稳定。求最终电压和热量。
\end{drill}
\fig{fig:em2:drill:4:s}{初始上极板尚未连通，分别带 $+CU$ 与 $-2CU$；闭合开关后经细电阻再分配}{\begin{tikzpicture}

\draw[world](0,0)--(4,0) (.5,0)--(.5,.8) (.5,1.2)--(.5,2.5)--(1.2,2.5) (2,2.5)--(2.5,2.5) (3.3,2.5)--(3.5,2.5)--(3.5,1.2) (3.5,.8)--(3.5,0);\draw[world](1.2,2.5)--(1.9,2.85);\draw[world](2.5,2.3)rectangle(3.3,2.7);\node[above]at(2.9,2.7){$r>0$};\foreach \xx/\lab in {.5/C,3.5/{2C}}{\draw[world]({\xx-.35},1.2)--({\xx+.35},1.2) ({\xx-.35},.8)--({\xx+.35},.8);\node[below]at(\xx,.55){$\lab$};}\node[left]at(.1,1.8){$+CU$};\node[right]at(3.9,1.8){$-2CU$};\node at(2,-.7){两下极板已连接；磁场与外源均不参与再分配};

\end{tikzpicture}}
\end{question}
\answerlink{em2:drill:4}
\begin{question}\begin{drill}[电阻改变后的圆环电势分配]\label{em2:drill:5}
半径 $a$ 的圆环处于轴对称向里增强的磁场中，总电动势 $e>0$。按逆时针，从 $P$ 到 $Q$ 的上半圆含两段四分之一圆，电阻分别为 $R,3R$；下半圆电阻为 $2R$。上端中点为 $M$。求 $V_P-V_M,V_M-V_Q,V_P-V_Q$，并计算电势最高的连接点。
\end{drill}
\fig{fig:em2:drill:5:s}{题设：电阻改变后的圆环电势分配}{\begin{tikzpicture}[scale=.88]\draw[world](0,0)circle(1.5);\fill(0,0)circle(1.3pt)node[below]{$O$};\draw[dim](0,0)--(1.5,0)node[midway,above]{$a$};\draw[flow](1.5*.707,1.5*.707)arc(45:110:1.5);\node[right]at(1.5,0){$P$};\node[left]at(-1.5,0){$Q$};\node[above]at(0,1.5){$M$};\node at(1.1,1.4){$R$};\node at(-1.1,1.4){$3R$};\node at(0,-2){$2R$};\end{tikzpicture}}
\end{question}
\answerlink{em2:drill:5}
\begin{question}\begin{drill}[由两次电压测量反求磁场变化率]\label{em2:drill:6}
圆形磁区内 $B(t)$ 均匀，半径 $a$。两个不改变网络负载的高内阻表连接相同端点，其两套固定引线之间围出面积 $S$，按第一路径前进、第二路径返回，法线向纸外。第一次有向读数 $U_1=0.80\,\mathrm V$，第二次 $U_2=0.50\,\mathrm V$，$S=0.015\,\mathrm{m^2}$，窄条完全位于磁区。求 $\dot B_z$；若另加一只包围全磁区的 $N$ 匝高阻探测线圈，求其电动势。
\end{drill}
\fig{fig:em2:drill:6:s}{题设：由两次电压测量反求磁场变化率}{\begin{tikzpicture}[scale=.88]\draw[world](0,0)node[left]{$A$}--(4,0)node[right]{$B$};\draw[guide](0,0)--(0,1.5)--(4,1.5)--(4,0);\node at(2,.8){$S,\ \odot z$};\node[below]at(2,0){路径1};\node[above]at(2,1.5){路径2};\end{tikzpicture}}
\end{question}
\answerlink{em2:drill:6}
\begin{question}\begin{drill}[感应驱动与直流电源相接]\label{em2:drill:7}
三支路并联于 $A,B$。第一支路为内阻 $R$ 的直流源，其沿 $A\to B$ 的电动势为 $E$；第二支路为电阻 $R$ 与同向感应电动势 $e$；第三支路为纯电阻 $R$。求三支路电流；当 $e$ 从零增大时，直流电源何时开始吸收电能？求对应临界状态下的总发热率。
\end{drill}
\fig{fig:em2:drill:7:s}{题设：感应驱动与直流电源相接}{\begin{tikzpicture}[scale=.88]\draw[world](0,0)--(4.8,0) (0,3.6)--(4.8,3.6);\node[above left]at(0,3.6){$A$};\node[below left]at(0,0){$B$};\draw[world](0.0,0)--(0.0,.8) (-0.18,.8)rectangle(0.18,1.6) (0.0,1.6)--(0.0,2.35) (0.0,3.05)--(0.0,3.6);\node[left]at(-0.25,1.2){$R$};\draw[world](0.0,2.7)circle(.35);\draw[flow](0.0,2.92)--(0.0,2.48);\node[right]at(0.45,2.7){$E$};\draw[flow](0.7,1.7)--(0.7,.6)node[midway,right]{$i_1$};\draw[world](2.4,0)--(2.4,.8) (2.2199999999999998,.8)rectangle(2.58,1.6) (2.4,1.6)--(2.4,2.35) (2.4,3.05)--(2.4,3.6);\node[left]at(2.15,1.2){$R$};\draw[world](2.4,2.7)circle(.35);\draw[flow](2.4,2.92)--(2.4,2.48);\node[right]at(2.85,2.7){$e$};\draw[flow](3.0999999999999996,1.7)--(3.0999999999999996,.6)node[midway,right]{$i_2$};\draw[world](4.8,0)--(4.8,.8) (4.62,.8)rectangle(4.9799999999999995,1.6) (4.8,1.6)--(4.8,2.35) (4.8,3.05)--(4.8,3.6);\node[left]at(4.55,1.2){$R$};\draw[world](4.8,2.7)circle(.35);\draw[flow](4.8,2.92)--(4.8,2.48);\node[right]at(5.25,2.7){$0$};\draw[flow](5.5,1.7)--(5.5,.6)node[midway,right]{$i_3$};\end{tikzpicture}}
\end{question}
\answerlink{em2:drill:7}
\begin{question}\begin{drill}[感应网络中的孤立中间导体]\label{em2:drill:8}
感应电阻网络已建立稳定端电势差 $U=V_A-V_B$。另一支路用串联电容 $C_1,C_2$ 接在 $A,B$，中间导体 $X$ 完全绝缘，其净电荷恒为 $Q_X$。沿该支路从 $A$ 到 $B$ 的全部固定引线感应积分为 $e_c$。电容电压 $u_1,u_2$ 均按 $A\to B$ 的电压降计。求两电压、储能，并讨论引线改变时的响应。
\end{drill}
\fig{fig:em2:drill:8:s}{题设：感应网络中的孤立中间导体}{\begin{tikzpicture}[scale=.88]\draw[world](0,0)--(1.7,0) (2.1,0)--(3.5,0) (3.9,0)--(5.5,0);\draw[world](1.7,-.35)--(1.7,.35) (2.1,-.35)--(2.1,.35) (3.5,-.35)--(3.5,.35) (3.9,-.35)--(3.9,.35);\node[above]at(1.9,.4){$C_1,u_1$};\node[above]at(3.7,.4){$C_2,u_2$};\node[left]at(0,0){$A$};\node[right]at(5.5,0){$B$};\fill(2.8,0)circle(.05)node[below=8pt]{$X,Q_X$};\node at(2.8,1.4){沿 $A\to X\to B$ 的引线积分为 $e_c$};\end{tikzpicture}}
\end{question}
\answerlink{em2:drill:8}
\begin{question}\begin{drill}[圆盘发电机的端压与电磁转矩]\label{em2:drill:9}
薄金属圆盘半径 $a$，在向纸外的匀强磁场 $B$ 中绕中心逆时针匀速转动，角速度 $\omega$。中心 $O$ 与圆周电刷 $A$ 经固定电阻 $R_L$ 闭合，盘内等效电阻为 $r$，接触电阻已计入，忽略自感。求端电压、电流和维持匀速所需转矩，并解释圆盘几何轮廓不变仍能发电的原因。
\end{drill}
\fig{fig:em2:drill:9:s}{题设：圆盘发电机的端压与电磁转矩}{\begin{tikzpicture}[scale=.88]\fill[gray!12](0,0)circle(1.6);\draw[world](0,0)circle(1.6);\draw[world](0,0)--(0,-2.4)--(3,-2.4)--(3,0)--(1.6,0);\fill(0,0)circle(.06)node[above left]{$O$};\fill(1.6,0)circle(.06)node[above right]{$A$};\node[draw,fill=white]at(3,-1.2){$R_L$};\draw[dim](0,0)--(1.6,0)node[midway,below]{$a$};\node at(-.8,.6){$\odot B$};\draw[flow](1.6,0)arc(0:70:1.6);\node at(1.2,1.8){$\omega$};\end{tikzpicture}}
\end{question}
\answerlink{em2:drill:9}
